\documentclass[a5paper,10pt]{article}

% https://github.com/InDevRus/filippov-solutions
% Нумерация страниц
\pagenumbering{gobble}

% Настройка полей
\usepackage{geometry}
\geometry{left=1cm}
\geometry{right=1cm}
\geometry{top=1cm}
\geometry{bottom=1.5cm}

% Вставка таблицы
\usepackage{array}
\usepackage{tabularx}

% Инструменты выравнивания формул
\usepackage{amsmath}
\usepackage{mathtools}

% Диакритика в математике
\usepackage{amsfonts}

% Вычисления размеров на ходу
\usepackage{calc}

% Удобные записи производных
\usepackage[italic=true]{derivative}

% Настройки сносок
\usepackage{enotez}

% Длинные знаки векторов
\usepackage{anyfontsize}
\usepackage[b]{esvect}

% Вставка картинок
\usepackage{graphicx}

% Вставка красной строки
\usepackage{indentfirst}
\setlength{\parindent}{1em}

% Условия в командах
\usepackage{xifthen}

% Луа-скрипты
\usepackage{luacode}

% Настройка команд отдельно в математическом и текстовом режимах
\usepackage{mathcommand}

% Для повторения LaTeX кода
\usepackage{multido}

% Конфигурация отступов формул
\delimitershortfall-1sp
\usepackage{mleftright}
\mleftright

% Растягивание объектов
\usepackage{relsize}

% Обтекание текстом
\usepackage{wrapfig}
\usepackage{subcaption}
\usepackage{floatflt}

% Поддержка ввода кириллицы
\usepackage[english, russian]{babel}

% Настройка корневой позиции для компилятора
\usepackage{import}
\providecommand*{\ProjectRootPath}{..}

\import{\ProjectRootPath/preambles/macros/}{theorems}

\import{\ProjectRootPath/preambles/fonts}{font_setup}

% Знак градуса
\usepackage{siunitx}

\import{\ProjectRootPath/preambles/macros/}{operators}
\import{\ProjectRootPath/preambles/macros/}{redefinitions}

\import{\ProjectRootPath/preambles/fonts}{letters_setup}
\import{\ProjectRootPath/preambles/fonts/adjustments}{custom_superscripts}

\import{\ProjectRootPath/preambles/custom}{formatting}
\import{\ProjectRootPath/preambles/custom}{frames}
\import{\ProjectRootPath/preambles/custom}{list_setup}

% TODO: перевести команды на современный стиль объявления

\begin{document}

Все рассуждения в решении этой задачи будут строиться в обозначениях, введённых в книге
Степанова В. В. <<Курс дифференциальных уравнений>>, Глава I, \S 2. 
Шаги доказательства также будут ссылаться на результаты этого параграфа.

\begin{framed}
    \begin{property}[интегрирующего множителя уравнения в дифференциальной форме]
        Если у уравнения
        $M\br{x; y} \d x + N\br{x; y} \d y = 0$
        есть интегрирующий множитель $\mu\br{x; y}$,
        непрерывный в выколотой окрестности некоторой точки $A$, а $U\br{x; y}$ -- интеграл 
        этого уравнения, то любой другой интегрирующий множитель $\sub{\mu}{1}\br{x; y}$ в этой же
        выколотой окрестности выражается формулой
        $$\sub{\mu}{1}\br{x; y} = f\br{U} \mu\br{x; y},$$
        где $f\br{U} = f\br{U\br{x; y}}$ -- произвольная функция, дифференцируемая всюду, кроме,
        возможно, точки $A$.
    \end{property}

    Доказательство этого факта можно найти в упомянутой выше книге, Глава I, \S 3.

    Более того, если множитель $\mu\br{x; y}$ обладает частными производными 
    $\parder{\mu}{x}$ и $\parder{\mu}{y}$, то функции $\sub{\mu}{1}\br{x; y}$, заданные равенством выше,
    суть общее решение уравнения в частных производных
    $$N \parder{\mu}{x} - M \parder{\mu}{y} = \br{\parder{M}{y} - \parder{N}{x}} \mu.$$
\end{framed}

\begin{framed}
    \begin{property}[сохранения интегрирующего множителя при замене переменных] \hspace{1mm}
        Если уравнение $M\br{x; y} \d x + N\br{x; y} \d y = 0$ имеет формальный непрерывный
        интегрирующий множитель
        $\mu\br{x; y}$,
        то непрерывно дифференцируемая замена переменных
        $x\br{\xi; \eta} = \f\br{\xi; \eta}$, $y\br{\xi; \eta} = g\br{\xi; \eta}$
        переводит исходное уравнение в такое, что
        $\mu\br{\f\br{\xi; \eta}; g\br{\xi; \eta}}$
        будет являться интегрирующим множителем и для него.
    \end{property}

    Действительно, если для некоторой функции $u\br{x; y}$ имеем 
    $du = d\br{u\br{x; y}} = \linebreak = M\mu \d x + N \mu \d y$,
    то в силу инвариантности формы первого дифференциала
    $$du = d\br{u\br{\f\br{\xi; \eta}; g\br{\xi; \eta}}} = 
    \br{M\, d\br{\f\br{\xi; \eta}} + N\, d\br{g\br{\xi; \eta}}} \cdot \mu.$$
\end{framed}
\pagebreak

В книге были проведены линейные замены переменных 
$\tsys{& \xi = \alpha x + \beta y \\& \eta = \gamma x + \delta y}$ 
с, вообще говоря, комплексными коэффициентами, сводящими уравнение к более простому виду. 
Такие замены геометрически представляют собой аффинные преобразования, 
которые только поворачивают и масштабируют исходную систему координат. 
При этом изначальная особая точка при таком преобразовании не меняет ни своё положение на 
координатной плоскости, ни свой тип.

Цель доказательства будет состоять в том, чтобы показать, что в случае особых точек
вида <<узел>>, <<фокус>>, <<вырожденный узел>> и <<дикритический узел>> наличие 
интегрирующего множителя, непрерывного в окрестности особой точки, приводит к противоречию. 
<<Центром>> особая точка быть не может (результат задачи \textbf{998}).

Во всех случаях далее обозначим непрерывный интегрирующий множитель исходного уравнения
за $\mu\br{x; y}$, $\sub{\lambda}{1}$ и $\sub{\lambda}{2}$ -- корни характеристического уравнения.

\begin{case}[1] Особая точка является <<узлом>>. \end{case}
В таком случае преобразование приводит к уравнению
$\dfrac {\d \eta} {\d \xi} = \dfrac {\sub{\lambda}{1} \eta} {\sub{\lambda}{2} \xi}$;
$\dfrac {\sub{\lambda}{1}} {\sub{\lambda}{2}} \eta \d \xi - \xi \d \eta = 0$.
Пусть $\dfrac {\sub{\lambda}{1}} {\sub{\lambda}{2}} = L > 0$.
Для уравнения
$- L \eta \d \xi + \xi \d \eta = 0$
интегрирующим множителем будет, например, $\sub{\mu}{1} = \tpow{\xi}{-L - 1}$, 
ибо после умножения будет
$-L \tpow{\xi}{-L - 1} \eta d \xi + \tpow{\xi}{L} \d \eta = 0$,
$d\br{\eta \tpow{\xi}{-L}} = 0$,
то есть общий интеграл равен $U\br{\xi; \eta} = \eta \tpow{\xi}{-L}$.

С другой стороны, $\mu\br{x; y} = \mu\br{x\br{\xi; \eta}; y\br{\xi; \eta}}$ -- также 
интегрирующий множитель, а значит для $f = f\br{\tau}$ имеем
$\sub{\mu}{1} = \tpow{\xi}{-L - 1} = f\br{\eta \tpow{\xi}{-L}} \mu\br{\xi; \eta}$.

Зафиксируем любое число $\tau \ne 0$, для которого значение $f\br{\tau}$ конечно и
рассмотрим \linebreak соответствующую непрерывную траекторию
$\eta = \tau \tpow{\xi}{L}$, $\eta \tpow{\xi}{-L} = \tau$.
Тогда в окрестности начала координат получаем противоречие
$\tpow{\xi}{-L - 1} = f\br{\tau} \mu\br{\xi; \tau \tpow{\xi}{L}}$,
так как правая часть полученного равенства непрерывна, а значит, ограничена, а левая является
неограниченной, ибо $-L - 1 < -1$.

В полученном противоречии заключается принципиальное различие между особыми 
точками вида <<узел>> и <<седло>>. Если особая точка является <<узлом>>, то все траектории 
проходят через особую точку, а через <<седло>> не проходит никакая, кроме двух сепаратрис.

Отдельно также отметим, что если особая точка является <<седлом>>, то непрерывный в окрестности
начала координат интегрирующий множитель существует всегда. В самом деле, без ограничения 
общности выберем $\sub{\lambda}1$ и $\sub{\lambda}2$ такими, что 
$\vbr{\dfrac {\sub{\lambda}{1}} {\sub{\lambda}{2}}} =
- \dfrac {\sub{\lambda}{1}} {\sub{\lambda}{2}} > 1$ (иначе поменяем их местами).
Тогда после преобразования имеем гарантированно непрерывный интегрирующий множитель
$$\sub{\mu}{1}\br{\xi; \eta} = \tpow{\xi}{-\frac {\sub{\lambda}{1}} {\sub{\lambda}{2}} - 1}
\text{, где }-\frac {\sub{\lambda}{1}} {\sub{\lambda}{2}} - 1 > 0.$$

\begin{case}[2]Особая точка является <<фокусом>>.\end{case}
В терминах, введённых в параграфе, комплексные корни обозначаются 
$\sub{\lambda}{1} = p + qi$ и $\sub{\lambda}{2} = p - qi$. 
Коэффициенты $\alpha$, $\beta$, $\gamma$ и $\delta$ получаются комплексными, 
поэтому вводится ещё одна замена переменных $\xi = u + iv$, $\eta = u - vi$, 
которая возвращает изначально вещественнозначные траектории в вещественную плоскость.

В результате имеем уравнение в новых терминах
$\br{qu - pv} du + \br{pu + qv} \d v = 0$.
С учётом того, что $q \ne 0$, обозначим $\dfrac {\hspace{0.20em}p} {q} = s$.
$\br{u - s v} du + \br{su + v} \d v = 0$.
Интегрирующим множителем этого уравнения является, например,
$$\sub{\mu}{1} = 2 \pow{e}{2s \arctg\br{\frac {v} {u}}},$$
потому что после домножения получаем полный дифференциал
\begin{multline*}
    \sub{\mu}{1} \cdot \br{u - sv} du + \sub{\mu}{1} \cdot \br{su + v} \d v = \\
    = 2 \pow{e}{2s \arctg\br{\frac {v} {u}}} \pow{e}{2s \arctg\br{\frac {v} {u}}} \br{u - sv} du
    + 2 \pow{e}{2s \arctg\br{\frac {v} {u}}} \br{su + v} \d v    = \\
    = d\ \br{\br{\pow{u}{2} + \pow{v}{2}} \pow{e}{2s \arctg\br{\frac {v} {u}}}} = dU.
\end{multline*}

Однако $\mu\br{u; v}$ -- также интегрирующий множитель, притом непрерывный, то есть
$$\mu\br{u; v} = f\br{U} \sub{\mu}{1} = f\br{\br{\pow{u}{2} + \pow{v}{2}}
\pow{e}{2s \arctg\br{\frac {v} {u}}}} \pow{e}{2s \arctg\br{\frac {v} {u}}}.$$

Выберем любую логарифмическую спираль 
$U = \br{\pow{u}{2} + \pow{v}{2}} \pow{e}{2s \arctg\br{\frac {v} {u}}} = 
\pow{\tau}{2}$, $\tau > 0$, 
для всех точек которой 
$\mu\br{u; v} = f\br{\pow{\tau}{2}} \pow{e}{2s \arctg\br{\frac {v} {u}}}$.
Всякая логарифмическая спираль обладает тем свойством, что она пересекается с прямой 
$v = ku$ бесконечно много раз в сколь угодно малой окрестности начала координат. 
Зафиксируем любое $k$ и рассмотрим последовательность таких пересечений 
$\sub{P}n = \br{\sub{u}n; \sub{v}n}$, сходящуюся к началу координат.
Итак, с одной стороны, $\sub{v}n = k \sub{u}n$ и $\tlim{n \to \infty} {\sub{P}n} = \br{0; 0}$,
а с другой для всех этих точек выполняется равенство
$\mu\br{\sub{u}n; \sub{v}n} = f\br{\pow{\tau}{2}} \pow{e}{2s \arctg k}$.
Если взять предел при $n \to \infty$, то
$\mu\br{0; 0} = \tlim{n \to \infty} \mu\br{\sub{u}n; \sub{v}n}
= f\br{\pow{\tau}{2}} \pow{e}{2s \arctg k}$. 
$\mu$ может быть непрерывна, только если правая часть не зависит ни от $\tau$, ни от $k$, 
то есть если $f\br{\pow{\tau}{2}} = 0$. Последнее означало бы, что $f\br{\tau}$ 
вместе с $\mu\br{u; v}$ -- тождественные нули, что невозможно.

\begin{case}[3]
    Особая точка является <<вырожденным узлом>>.
\end{case}

Преобразование координат из $\br{x; y}$ в $\br{\xi; \eta}$ приводит уравнение к следующему виду: 
$\br{\xi + \lambda \eta} \d \xi - \lambda \xi \d \eta = 0$,
где $\lambda = \sub{\lambda}{1} = \sub{\lambda}{2}$, $\lambda \ne 0$ и $\lambda \ne 1$.

Для получения противоречия можно взять любой из интегрирующих множителей последнего уравнения,
например $\sub{\mu}{1} = \dfrac {1} {\pow{\xi}{2}}$ или $\sub{\mu}{2} = \exp\br{-2\lambda \dfrac {\eta} {\xi}}$. Рассмотрим, например, $\sub{\mu}{1}$.

В параграфе показано, что общим интегралом этого уравнения будет 
$U\br{\xi; \eta} = \linebreak = \dfrac {1}{\lambda} \ln \vbr{\xi} - \dfrac {\eta}{\xi}$.
Таким образом, $\mu\br{\xi; \eta} 
= f\br{\dfrac {1}{\lambda} \ln \vbr{\xi} - \dfrac {\eta}{\xi}} \dfrac {1} {\pow{\xi}{2}}$.

Все траектории в этом случае проходят через особую точку, так что точки на кривой 
$\dfrac {1}{\lambda} \ln \vbr{\xi} - \dfrac {\eta}{\xi} = \tau$ с зафиксированным
$\tau \ne 0$ в окрестности начала координат дают
$\mu\br{\xi; \eta} = \mu\br{\xi; \dfrac {1} {\lambda} \xi \ln\vbr{\xi} - \tau \xi} 
= f\br{\tau} \dfrac {1} {\pow{\xi}{2}}$.
Последнее влечёт противоречие, так как левая часть непрерывна в окрестности начала координат,
а правая -- неограничена.
 
\begin{case}[4]
    Особая точка является <<дикритическим узлом>>.
\end{case}

Имеем $\sub{\lambda}{1} = \sub{\lambda}{2} = 1$ и исходное уравнение $y \d x - x \d y = 0$.
В качестве интегрирующего множителя можно рассмотреть 
$\sub{\mu}{1} = \dfrac {1} {\pow{x}{2}}$, $\sub{\mu}{2} = \dfrac {1} {\pow{y}{2}}$ 
или $\sub{\mu}{3} = \dfrac {1} {xy}$. При этом $U\br{x; y} = \dfrac {\hspace{0.1em}y} {x}$.

Например, вдоль прямых $y = \tau x$ получаем противоречие
$\mu\br{x; y} = f\br{\tau} \dfrac {1} {\pow{x}{2}}$,
аналогичное тому, что возникло в случае <<узла>>.

Поскольку кроме рассмотренных видов особых точек никаких других быть не может,
утверждение доказано.

\end{document}
